\section{10月13日勘误与证据范围}\label{sec:reader-v2-errata}
新版正文回应总体、全50表和六地图审查。10月9日公开PDF及77页中文补充原件
继续保留于\path{revision46-review-v1}，更正PDF另行版本化；
不把历史核对条目改写成当时已经包含正确结论。

上一版正文ES分母$2N^2$及对应检查条目有误。执行内容\texttt{c3e44393}开头的绑定、
保留运行设置与公开配对复算收据共同绑定\path{metrics.py}至
\texttt{c00059e4}开头的SHA-256；共同评分器\path{formal_scoring.py}
也匹配执行登记的\texttt{d39adc6d}开头哈希。
其调用链采用有序粒子对U统计量，分母$2N(N-1)$，与原科学补充一致；
惯性点质量评分另用直接距离分支。完整哈希及两个纯函数源码见
\path{reader-v2-score-binding.json}。
本次是源码绑定与文稿勘误，不是新原始数组审计或重算科学评分。
新增两粒子解析样例能区分此算子与经验分布估计式。

G01改为实际建筑／植被／水体湿地／有效其他分组，不再称已画出坡度大小。
正文补入输入几何、十配置设计、全部21评分／位置取舍、校准／大小与种子／块不确定性例子。
热力图保留原值和角色排序，统一浅色、黑字及对称的带符号覆盖偏差色阶。
表题就地分开历史待核状态与已完成标量／审计收据，未重建缺失的第二bootstrap流细节。

R05代理与标签错配、R06惩罚／表示混杂、R10缺少任务匹配概率基线，
在\path{revision46-response-v3.json}中明确保留为科学缺口，
分别对应既有N03、N04、N01后续设计；本次未执行这些实验。
六幅私有地图保留原案例、目标时刻、每面板512粒子及漏覆盖例子，
本地版改用共享图例、可读标签和不同路线／预测圆线型。
地图、坐标与私有清单继续留在公开包之外。
真实署名、外部读者反馈、路线权限和人工验收仍待确认。
